\scanpage{12}
\uline{NB} \deleted{\deleted{\unsure{Soit}} \unsure{Supp.|Supposons} $S=\operatorname{Spec}k$, \deleted{\unsure{$X=$}} $k$ corps alg. \uline{clos}, et $X=S$,}
\[
\begin{aligned}
\deleted{\unsure{S_0}}&\quad\deleted{S_0=\text{schémas \inserted{de t.f.} étales sur }X=\text{\unsure{isom.} locaux sur }X,}\\
&\quad\deleted{\mathcal S=\text{schémas de t.f. sur }S}\\
&\quad\deleted{\mathcal M=\text{imm. ouvertes}}
\end{aligned}
\]

\deleted{La catégorie $\mathcal C$ est celle des schémas \deleted{\unsure{ét.}} finis, sur $k$, avec la top. de Zariski, un faisceau \inserted{de $\mathcal O$-modules} dessus revient à la donnée d'un \deleted{module} foncteur $\Lambda\leadsto F(\Lambda)$, $\Lambda$ fini sur $k$, commutant aux produits.}

\deleted{Les faisceaux qu. coh. \inserted{sur $\operatorname{Spec}\Lambda$} sont ceux \unsure{pour} \deleted{\unsure{lesquels}} de la forme $\deleted{\unsure{F}}\ \Lambda'\leadsto M\otimes_{\Lambda}\Lambda'$, avec $M$ un $\Lambda$-module, les faisceaux de prés. finie sont ceux pour lesquels $\deleted{\dim_k}M$ \deleted{\unsure{finie}} \inserted{de t.f.}. Mais $\mathcal O$ \unsure{lui-m.} n'est pas cohérent, car \deleted{\illegible{1}} soit un hom $\mathcal O_\Lambda^n\overset u{-}\mathcal O_\Lambda$ défini par \deleted{\illegible{1}} $u_1\ldots u_n\in\Lambda$, alors $\operatorname{Ker}u$ est un ``faisceau sur \deleted{Spec} $\Lambda$'' \unsure{f.|?}}
\[
\deleted{\Lambda'\longmapsto\operatorname{Ker}\bigl(u_{\Lambda'}:\deleted{\unsure{\mathcal O_\Lambda}}\Lambda'{}^n-\Lambda'\bigr)}
\]
\deleted{qui n'est pas en général déduit de $\operatorname{Ker}u$ comme $\Lambda'\leadsto(\operatorname{Ker}u)\otimes_\Lambda\Lambda'$, manque de platitude \ldots\,.}

% margin note written sideways
\margin{\diagram{\begin{tikzcd}[ampersand replacement=\&,column sep=small]
T'_y\arrow[d,dash] \& T'\arrow[d,"\unsure{\text{ét.}}"'] \& X'\arrow[d] \& x'\\
T_y \& T \& X \& x
\end{tikzcd}}
$U\cap T_y$ dense dans $T_y$\\
$U'\cap T'_y$ dense \unsure{ds|dans} $T'_y$\\
$T'_{\unsure{k|i}}=\bigcup T'_i$\\
$T_y=\bigcup T'_{iy}$}
% margin note written sideways, in the opposite direction to the main text
\margin{\diagram{\begin{tikzcd}[ampersand replacement=\&,column sep=small]
T \&[-1em] \&\\
\tikz[baseline=(a.base)]{\node[inner sep=1pt](a){$\overline{\{x'\}}$};\draw (a.center) ellipse (1.0em and 1.3em);}\cap X_y \& U_y \&\\[-.5em]
\shortparallel \& \&\\[-.5em]
T_y \& \& T_{y\unsure{a}}
\end{tikzcd}}}
% archive sheet number: 11
